\subsection{Models}

We compare two models from the same architecture family:
\begin{itemize}
  \item \textbf{Gemma-3-4B-IT}~\citep{gemma2025}: a 4B-parameter general-purpose
        instruction-tuned VLM.
  \item \textbf{MedGemma-4B-IT}~\citep{medgemma2025}: a medical fine-tune of
        Gemma-3-4B trained on medical imaging and clinical text.
\end{itemize}
Because MedGemma-4B was fine-tuned \emph{directly} from Gemma-3-4B, the pair shares
architecture, layer count (34 decoder layers), tokenizer, and image token budget
(256 visual tokens per image), making the information-flow curves directly comparable.

\textbf{Note on encoder difference.} MedGemma uses a medical SigLIP vision encoder,
while Gemma-3 uses a general SigLIP encoder. The probe measures routing through the
\emph{decoder}, so the reported information-flow values reflect ``how the decoder uses
the image features it receives,'' not properties of the encoder itself.

\subsection{Dataset}

We evaluate on \textbf{VQA-RAD}~\citep{lau2018dataset}, a radiology visual
question-answering dataset covering chest X-rays, CT scans, and MRIs. We use 250
samples from the VQA-RAD test split, partitioned by the written protocol into a
100-sample discovery split (indices 0000--0099) and a 150-sample confirmation split
(indices 0100--0249), with identical indices for both models. Accuracy is computed on
the closed-ended yes/no questions only.

\subsection{Setup}

Both models are run in \texttt{bfloat16} (no weight quantization) on Kaggle T4 GPUs.
The forward hook is applied with window size $K=1$, sweeping all 34 decoder layers
per sample per block type. Generation uses greedy decoding with
\texttt{max\_new\_tokens=64}.

\subsection{Pre-specified analysis protocol}
\label{sec:hypotheses}

\textbf{Timing of the protocol.} The work ran in four steps. (1) A 50-sample pilot of the probe on both models. It measured closed-question accuracy,
F1 and QDR. Accuracy and F1 were computed with a text-extraction scorer, which was later
found to be invalid; the paper reports only the logit verdict. The pilot's QDR drop was
larger than the effects later measured, and it did not replicate at $n{=}100$. QDR was then
chosen as the primary metric. (2) The full pipeline was rerun on discovery indices 0000--0099 (100
samples per model), which produced the candidate metrics. (3) The protocol
(\texttt{PREREGISTRATION.md}) was written on 2026-10-24, after the Gemma 100-sample trace
had been analysed, the 50-sample MedGemma trace and the smoke-test outputs had been seen, and the MedGemma 100-sample output had not yet been analysed. (4) A launch on
0100--0199 was started and stopped after about three minutes, before the model loaded, and
produced no output. Confirmation then ran once on 0100--0249 (150 samples, none used in
discovery). The protocol has not been filed with an external registry, so the confirmation split is the clean test of the
pre-specified claims.

\textbf{Deviation from protocol.} Two of the eleven declared metrics turned out to be
deterministic transforms of two others (answer perplexity of mean answer log-prob;
image entry depth of image entry layer), carrying identical test statistics. We
therefore apply the Holm correction across the nine distinct quantities rather than the
eleven declared metrics. Under the original eleven-test correction the same five effects
survive, so the change does not alter any conclusion.

Our analysis follows a written protocol dated 2026-10-24, which was written after part of the discovery analysis (see the timing note above). It has not been filed with an external registry. The VQA-RAD test split is partitioned into a
\textbf{discovery} set (indices 0000--0099), a \textbf{confirmation} set
(0100--0249), and a reserve (0250--0450, not used). All metrics are computed and tested on discovery; any
metric with a Holm-adjusted $p<0.05$ across the nine distinct quantities (see the deviation below) becomes a
\emph{discovery candidate} that is re-tested, with the same test and direction, on the
confirmation set, and is reported as a \emph{finding} only if it replicates
($p<0.05$, same direction). The pre-specified \textbf{primary} confirmatory test is
that the per-sample question dominance ratio (QDR, \texttt{image\_to\_last}) is lower
for MedGemma than for Gemma (Mann--Whitney $U$, two-sided, $\alpha{=}0.05$). The
pre-specified protocol also asks for a permutation test alongside each Mann--Whitney test; we
use a two-sided permutation test on the difference of medians (10{,}000 permutations) as
that companion check, and report it next to each raw $p$-value.

We state three directional expectations about what medical fine-tuning should change
(reported against, not assumed). They appear in the paper by
2026-10-26, but they are not in the protocol file as it stood that day. They are recorded in
the protocol by Amendment 4 (2026-10-29). The repository does not establish when they were
first written, so they are not claimed as pre-specified:
\begin{itemize}
  \item \textbf{H1} (Image Entry Layer): fine-tuning shifts image evidence to
        \emph{earlier} decoder depths (faster visual integration).
  \item \textbf{H2} (Question Dominance): fine-tuning lowers the QDR, making the
        model more visually grounded relative to language priors.
  \item \textbf{H3} (Peak Layer): the peak image-information layer shifts earlier.
\end{itemize}

\subsection{Accuracy Evaluation}

For closed-ended questions, accuracy uses the logit verdict: the model's probability of
``yes'' versus ``no'' at the first answer token, aggregated over casing and spacing
variants. Text extraction is not used for accuracy, because generations are truncated
before the answer. Open-ended questions are not evaluated. The routing metrics operate on
token log-probabilities and are unaffected by this choice.
